%%%%%%%%%%%%%%%%%
% This is an sample CV template created using altacv.cls
% (v1.1.4, 27 July 2018) written by LianTze Lim (liantze@gmail.com). Now compiles with pdfLaTeX, XeLaTeX and LuaLaTeX.
% 
%% It may be distributed and/or modified under the
%% conditions of the LaTeX Project Public License, either version 1.3
%% of this license or (at your option) any later version.
%% The latest version of this license is in
%%    http://www.latex-project.org/lppl.txt
%% and version 1.3 or later is part of all distributions of LaTeX
%% version 2003/12/01 or later.
%%%%%%%%%%%%%%%%
%% If you need to pass whatever options to xcolor
\PassOptionsToPackage{dvipsnames}{xcolor}

%% If you are using \orcid or academicons
%% icons, make sure you have the academicons 
%% option here, and compile with XeLaTeX
%% or LuaLaTeX.
% \documentclass[10pt,a4paper,academicons]{altacv}

%% Use the "normalphoto" option if you want a normal photo instead of cropped to a circle
% \documentclass[10pt,a4paper,normalphoto]{altacv}

\documentclass[10pt,a4paper]{altacv}
%% AltaCV uses the fontawesome and academicon fonts
%% and packages. 
%% See texdoc.net/pkg/fontawecome and http://texdoc.net/pkg/academicons for full list of symbols.
%% 
%% Compile with LuaLaTeX for best results. If you
%% want to use XeLaTeX, you may need to install
%% Academicons.ttf in your operating system's font 
%% folder.

\renewcommand{\emailsymbol}{\faIcon{at}}


% Change the page layout if you need to
\geometry{left=1cm,right=8.1cm,marginparwidth=6.1cm,marginparsep=1.2cm,top=1.25cm,bottom=1.25cm,footskip=2\baselineskip}

% Change the font if you want to.

% If using pdflatex:
%\usepackage[T1]{fontenc}
%\usepackage[utf8]{inputenc}
%\usepackage[default]{lato}
%\usepackage{hyperref} 
%\usepackage{url}
% If using xelatex or lualatex:
\setmainfont{Lato}

% Change the colours if you want to
\definecolor{Mulberry}{HTML}{72243D}
\definecolor{SlateGrey}{HTML}{2E2E2E}
\definecolor{LightGrey}{HTML}{666666}
\colorlet{heading}{Sepia}
\colorlet{accent}{Mulberry}
\colorlet{emphasis}{SlateGrey}
\colorlet{body}{LightGrey}

% Change the bullets for itemize and rating marker
% for \cvskill if you want to
\renewcommand{\itemmarker}{{\small\textbullet}}
\renewcommand{\ratingmarker}{\faCircle}
\newcommand\rurl[1]{%
  \href{https://#1}{\nolinkurl{#1}}%
}

\addbibresource{sample.bib}
\definecolor{mediumturquoise}{rgb}{0.28, 0.82, 0.8}
\definecolor{cyan}{rgb}{0.0, 1.0, 1.0}
\usepackage{xcolor}
\hypersetup{
    colorlinks,
    linkcolor={red!50!black},
    citecolor={blue!50!black},
    urlcolor= {blue!70!black} %{blue!80!black}
}

\usepackage{url}

%Annotate the necesssary ones in the bibtext directly from here: https://tex.stackexchange.com/questions/73136/make-specific-author-bold-using-biblatex/304968
\renewcommand*{\mkbibnamegiven}[1]{%
  \ifitemannotation{highlight}
    {\textbf{#1}}
    {#1}}

\renewcommand*{\mkbibnamefamily}[1]{%
  \ifitemannotation{highlight}
    {\textbf{#1}}
    {#1}}


\begin{document}
%


\name{Nikolay Bogoychev}
%\tagline{PhD student looking for post doctoral position}
\tagline{Research Interest: Machine Translation, GPGPU, GEMM, Low/high level performance optimisation }
%\photo{2.8cm}{Globe_High}
\personalinfo{%
  % Not all of these are required!
  % You can add your own with \printinfo{symbol}{detail}
  \email{\href{mailto:nheart@gmail.com}{nheart@gmail.com}}
  \location{United Kingdom}
  \homepage{\url{www.nbogoychev.com}}
  \github{\href{https://github.com/XapaJIaMnu}{XapaJIaMnu}}
  \stackoverflow{\href{https://stackoverflow.com/users/1846228/xapajiamnu}{XapaJIaMnu}}
  \threads{\href{https://threads.net/@xapajiamnu}{@xapajiamnu}}
  \bluesky{\href{https://bsky.app/profile/xapajiamnu.bsky.social}{@xapajiamnu.bsky.social}}
  \mastodon{\href{https://social.linux.pizza/@xapajiamnu}{@xapajiamnu@social.linux.pizza}}
  \scholar{\href{https://scholar.google.com/citations?user=z2hWTDMAAAAJ}{Nikolay Bogoychev}}
  \linkedin{\href{https://www.linkedin.com/in/nikolay-bogoychev-01831023}{nikolay-bogoychev}}
  %% You MUST add the academicons option to \documentclass, then compile with LuaLaTeX or XeLaTeX, if you want to use \orcid or other academicons commands.
%   \orcid{orcid.org/0000-0000-0000-0000}
}

%% Make the header extend all the way to the right, if you want. 
\begin{fullwidth}
\makecvheader
\end{fullwidth}

%% Depending on your tastes, you may want to make fonts of itemize environments slightly smaller
% \AtBeginEnvironment{itemize}{\small}


%% Provide the file name containing the sidebar contents as an optional parameter to \cvsection.
%% You can always just use \marginpar{...} if you do
%% not need to align the top of the contents to any
%% \cvsection title in the "main" bar.
\cvsection[page1sidebar]{Education}

\cvevent{PhD}{School of Informatics, University of Edinburgh}{2015 -- 2019}{Edinburgh,United Kingdom}
\begin{itemize}
\item Thesis: Fast machine translation on parallel and massively parallel hardware

\item Supervisors: Adam Lopez, Kenneth Heafield
\end{itemize}

%\divider

\cvevent{BSc}{Artificial Intelligence and Computer science, University of Edinburgh }{2010 --  2014}{Edinburgh,United Kingdom}
\begin{itemize}
\item First Class Honours
\end{itemize}

%\divider

\cvsection{Work Experience}

\cvevent{CTO}{Last Token}{June 2026 - Present}{Remote, United Kingdom}
Providing premium quality training data for Large language models.

\divider

\cvevent{Research Scientist}{Meta}{December 2023 - December 2025}{London, United Kingdom}
Working on the LLaMa family of models.
\begin{itemize}
    \item Released the LLaMa 3.2 1B and 3B models
    \item Created the \href{https://www.kaggle.com/datasets/metaresearch/multiloko}{\textit{Multiloko}} multilingual local knowledge evaluation dataset
    \item Also worked on Pruning, Distillation, Evaluation, Scaling Laws, etc
\end{itemize}

\divider

\cvevent{CTO}{Efficient Translation Limited}{March 2023 - December 2023}{Edinburgh, United Kingdom}
Edge device machine translation for Privacy focussed customers.
\begin{itemize}
    \item x86/arm fast machine translation services deployed on local devices.
    \item terminology translation.
\end{itemize}

\divider

\cvevent{Postdoctoral Researcher}{University of Edinburgh, School of Informatics}{January 2019 - December 2023}{Edinburgh, United Kingdom}
Working on neural machine translation and High performance computing
\begin{itemize}
    \item translateLocally, cross platform, fast, offline desktop translation software {\rurl{github.com/XapaJIaMnu/translateLocally}}
    \item marian NMT 8bit integer CPU/GPU tensorcore inference code
    \item CPU int8 VNNI GEMM: {\rurl{github.com/kpu/intgemm}}
    \item Experimental GEMM tiling: {\rurl{github.com/XapaJIaMnu/bfTile}}
    \item Bergamot: Privacy driven Firefox translation. For details: {\rurl{browser.mt}}
\end{itemize}

\divider

\cvevent{Consultant}{Intel}{March 2019 - December 2023}{Edinburgh, United Kingdom}
Exploring opportunities for CPU training and inference for NMT:
\begin{itemize}
    \item BF16 and other quantisation schemes for training and inference.
    %\item Profile, optimise and vectorise NMT codebase.
    \item Low level hardware specific optimisation with focus on NMT.
\end{itemize}


%\medskip

\cvsection{Selected Publications}
%\nocite{me_sgd}
\begin{itemize}
\item \fullcite{hupkes2025multilokomultilinguallocalknowledge}
\item \fullcite{dubey2024llama3herdmodels}
\item \fullcite{bogoychev-etal-2024-ups}
\item \fullcite{bogoychev2023opuscleaneropustraineropensource}
\item \fullcite{grivas-etal-2022-low}
\item \fullcite{bogoychev-2021-parameters}
\item \fullcite{bogoychev-etal-2021-translatelocally}
\item \fullcite{bogoychev-chen-2021-highs}
\item \fullcite{aji-etal-2020-neural}
\item \fullcite{bogoychev-etal-2020-edinburghs}
\item \fullcite{me_sgd}
\item \fullcite{DBLP:conf/acl/BogoychevL16}
\end{itemize}
%\nocite{aji-etal-2020-neural}
%\nocite{DBLP:conf/acl/BogoychevL16}
%\begin{enumerate}
%    \item \bibentry{me_sgd}
%    \item \bibentry{aji-etal-2020-neural}
%    \item \bibentry{DBLP:conf/acl/BogoychevL16}
%  \end{enumerate}

\cvsection[page2sidebar]{Talks and Tutorials}

\cvevent{The Role of Locally Relevant Content in Evaluating Multilingual LLMs}{LocWorld54}{09/2025}{Monterey, United States}
{\rurl{https://nbogoychev.com/files/publications/locworld_multilko.pdf}}

\divider

\cvevent{Efficient Machine translation tutorial}{MT Marathon}{11/2021}{Edinburgh, United Kingdom}
{\rurl{nbogoychev.com/efficient-machine-translation/}}


%\cvsection{Other Experience}
%\cvevent{Consultant}{myLanguage}{12/2020}{Edinburgh, United Kingdom}
%\begin{itemize}
%    \item Porting x86\_64 software to ARM.
%    \item Providing knowledge about neural machine translation model training.
%\end{itemize}

\divider

%\cvevent{Parallelism Research intern}{Intel}{06/06/2017 - 06/10/2017}{Santa Clara, California, USA}
%Worked on Neural machine translation training on CPUs. Produced a detailed breakdown and the communication and computation costs of each operation in neural machine translation training. Developed methods for reducing internode communication during training.

%\divider

%\cvevent{Machine Learning Intern}{Amazon}{01/06/2016 - 30/09/2016}{Berlin, Germany}
%Optimising Spark/Scala Bayesian inference framework.

%\divider

%\cvevent{Research Assistant}{University of Edinburgh}{01/06/2014 - 01/12/2016}{Edinburgh, United Kingdom}
%Working on Statistical Machine translation
%\begin{itemize}
%    \item Multithreading optimisation for Moses.
%    \item Bilingual language model feature function for Moses.
%\end{itemize}

%\divider

%\cvevent{Intern}{Contemplate}{01/06/2013 - 01/09/2013}{Edinburgh, United Kingdom}
%Developing a JS frontend for a Java static analysis toolkit.

%\divider

%\cvevent{Language Intern}{SwiftKey}{01/06/2012 - 01/09/2012}{London, United Kingdom}
%Working on predictive typing framework for mobile phones.
%\begin{itemize}
%    \item Create language models for various language pairs.
%    \item Developing intelligent web crawler in clojure.
%\end{itemize}

\nocite{*}



%\printbibliography[heading=none,title={},maxnames=2]%,type=inproceedings]


%% If the NEXT page doesn't start with a \cvsection but you'd
%% still like to add a sidebar, then use this command on THIS
%% page to add it. The optional argument lets you pull up the 
%% sidebar a bit so that it looks aligned with the top of the
%% main column.
% \addnextpagesidebar[-1ex]{page3sidebar}


\end{document}
